\documentclass[11pt]{article}
\usepackage{amsfonts}
\usepackage[T1]{fontenc}
\usepackage{mathabx,graphicx}

\newcommand{\test}{\circlearrowright}
\def \loop {\ensuremath{\rotatebox[origin=c]{-90}{$\circlearrowright$}}}
\def \nestedloop {\ensuremath{\rotatebox[origin=c]{-90}{$\circlearrowright$}}^n}

\begin{document}

\section*{Parallel Computation}






% Complexity; Time Complexity; Space Complexity; Total Complexity Dimension > 1
\section{Complexity}

\subsection{Time Complexity of a Decision Problem $O_T \lbrack n \rbrack$}
Define Time Complexity $O_T [n]$ of solution $s^+$ to Decision Problem $D$ as the total number of logical operations
\begin{center}
\vspace{1mm}
$
X_i = \{x_1,...,x_n,C\}; \hspace{2mm} \hat{X}_i = \{x_1,...,x_{n+1},C\}
$
\\ \vspace{4mm}
$
D := f \lbrack X_i \rbrack \rightarrow a_{o} \in \{\mathbb{T}, \mathbb{F}\} \hspace{3mm} \forall X_i
$
\\ \vspace{4mm}
$
s^+[X_i] := P :
$
\\ \vspace{2mm}
$
(P \lbrack X_i \rbrack \rightarrow y_{o} == a_{o} \hspace{3mm} \forall X_i) \hspace{2mm} \cap \hspace{2mm} (P[\hat{X}_i] \supseteq P[X_i] \hspace{3mm} \forall X_i,\hat{X}_i)
$
\\ \vspace{4mm}
$
s^+ = \{ s_1,s_2,...,s_N|b_1,b_2,...,b_M,y_o\} = \{ s_1,s_2,...,s_{O_T \lbrack n \rbrack }, b_1, b_2,...,b_{O_S \lbrack n \rbrack},y_o \}
$
\\ \vspace{2mm}
$
= \{ \mathcal{L},\mathcal{M},y_o\}
$
\\ \vspace{3mm}
$
O_T[n] := |\mathcal{L}| = N
$
\end{center}

\subsection{Space Complexity $O_S \lbrack n \rbrack$}
Define Space Complexity $O_S \lbrack n \rbrack$ of solution $s^+$ to Decision Problem $D$ as the total number of memory elements
\begin{center}
\vspace{1mm}
$
X_i = \{x_1,...,x_n,C\}; \hspace{2mm} \hat{X}_i = \{x_1,...,x_{n+1},C\}
$
\\ \vspace{4mm}
$
D := f \lbrack X_i \rbrack \rightarrow a_{o} \in \{\mathbb{T}, \mathbb{F}\} \hspace{3mm} \forall X_i
$
\\ \vspace{4mm}
$
s^+[X_i] := P :
$
\\ \vspace{2mm}
$
(P \lbrack X_i \rbrack \rightarrow y_{o} == a_{o} \hspace{3mm} \forall X_i) \hspace{2mm} \cap \hspace{2mm} (P[\hat{X}_i] \supseteq P[X_i] \hspace{3mm} \forall X_i,\hat{X}_i)
$
\\ \vspace{4mm}
$
s^+ = \{ s_1,s_2,...,s_N|b_1,b_2,...,b_M,y_o\} = \{ s_1,s_2,...,s_{O_T \lbrack n \rbrack }, b_1, b_2,...,b_{O_S \lbrack n \rbrack},y_o \}
$
\\ \vspace{2mm}
$
= \{ \mathcal{L},\mathcal{M},y_o\}
$
\\ \vspace{3mm}
$
O_S[n] := |\mathcal{M}|+|y_o|^* = M + 1
$
\end{center}
\vspace{2mm}
*It is convention to reserve one memory element for output $y_o$.\\
Void programs do not require the $y_o$ memory element for output




\section{Definition of Complexity}
Define Complexity $O[n]$ as a vector of dimension V
\begin{center}
$
\bold{O}[n] := \hspace{3mm} < O_T [n], O_S [n],O_3[n],O_4[n]...,O_V[n]>
$
\end{center}

%\section{Total Complexity}
%\begin{center}
%$O[n] := O_T[n] + O_S[n] + \sum_{i=3}^{V} O_i[n]$
%\end{center}





















% Simple Computational Complexity; Restate Time and Space Complexity; (Simple) Total Complexity
\newpage
\section{Simple Computational Complexity}
The remainder of this document assumes simple computational complexity of dimension 2

\subsection{Definition}
Define simple computational complexity of dimension 2
\begin{center}
$
\bold{O}[n] := \hspace{3mm} < O_T [n], O_S [n] >
$
\end{center}




\subsection{Total Complexity}
Define Total Complexity of solution $s^+$
\begin{center}
$O[n] :=  |s^+[X_n]| = |\{\mathcal{L},\mathcal{M},y_o\}|$
\\ \vspace{4mm}
$
= |\mathcal{L}| + |\mathcal{M}| + |y_o| = N + M + 1
$
\end{center}






\subsection{Time Complexity}
Restate definition of Time Complexity $O_T[n]$ of solution $s^+$
\begin{center}
$
s^+ = \{ \mathcal{L},\mathcal{M},y_o\}
$
\\ \vspace{3mm}
$
O_T[n] := |\mathcal{L}| = N
$
\end{center}



\subsection{Space Complexity}
Restate definition of Time Complexity $O_S[n]$ of solution $s^+$
\begin{center}
$
s^+ = \{ \mathcal{L},\mathcal{M},y_o\}
$
\\ \vspace{2mm}
$
O_S[n] := |\mathcal{M}| + |y_o| = M + 1
$
\end{center}



\subsection{Total Complexity as a Function of Time and Space Complexity}
\begin{center}
\vspace{4mm}
$
O[n] :=  |s^+[X_n]| = |\{\mathcal{L},\mathcal{M},y_o\}|
$
\\ \vspace{2mm}
$
= |\mathcal{L}| + |\mathcal{M}| + |y_o|
$
\\ \vspace{2mm}
$
= O_T[n] + O_S[n]
$
\end{center}


\subsection{$O_S[n] > 0^*$}
Assuming Program is not void
\subsubsection{Proof}
Assume $O_S[n]$ = 0
\begin{center}
$
O_S[n] := |\mathcal{M}| + |y_o| 
$
\\ \vspace{2mm}
$
O_S[n] = 0 \Rightarrow \mathcal{M} = y_o = \emptyset
$
\\ \vspace{2mm}
$
y_o = \emptyset; \hspace{2mm} y_o \in \{\mathbb{T},\mathbb{F}\}$ by definition of $s^+$
\\ \vspace{4mm}
$
\therefore O_S[n] = 0$ contradicts the definition of solution $s^+$ of a decision problem
\\ \vspace{2mm}
$
O_S[n] \geq 0$ by definition of magnitude
\\ \vspace{2mm}
$
\therefore O_S[n] > 0
$
\end{center}


\subsection{$O_T[n] > 0^*$}
Assuming Program is not void
\subsubsection{Proof}
Assume $O_T[n] = 0$
\begin{center}
$
O_T[n] := |\mathcal{L}| 
$
\\ \vspace{2mm}
$
O_T[n] = 0 \Rightarrow y_o \not \in \{\mathbb{T},\mathbb{F}\}
$
\\ \vspace{2mm}
$
y_o \not \in \{\mathbb{T},\mathbb{F}\}; \hspace{2mm} y_o \in \{\mathbb{T},\mathbb{F}\}$ by definition of $s^+$
\\ \vspace{4mm}
$
\therefore O_T[n] = 0$ contradicts the definition of solution $s^+$ of a decision problem
\\ \vspace{2mm}
$
O_T[n] \geq 0$ by definition of magnitude
\\ \vspace{2mm}
$
\therefore O_T[n] > 0
$
\end{center}




\subsection{$O[n] > 0^*$}
Assuming Program is not void
\subsubsection{Proof}
\begin{center}
$
O[n] := O_T[n] + O_S[n]
$
\\ \vspace{3mm}
$
O_T[n] > 0; \hspace{2mm} O_S[n] > 0
$
\\ \vspace{3mm}
$
\therefore O[n] > 0
$
\end{center}

\subsection{$O[n] > O_T[n]*$}
Assuming Program is not void
\subsubsection{Proof}
\begin{center}
$
O[n] := O_T[n] + O_S[n]
$
\\ \vspace{3mm}
$
\hspace{2mm} O_S[n] > 0
$
\\ \vspace{3mm}
$
\therefore O[n] > O_T[n]
$
\end{center}

\subsection{$O[n] > O_S[n]^*$}
Assuming Program is not void
\subsubsection{Proof}
\begin{center}
$
O[n] := O_T[n] + O_S[n]
$
\\ \vspace{3mm}
$
O_T[n] > 0
$
\\ \vspace{3mm}
$
\therefore O[n] > O_S[n]
$
\end{center}


\subsection{$O[n+1] \geq O[n]$}
\subsubsection{Proof}
\begin{center}
$
X_i = \{x_1,...,x_n,C\}; \hspace{2mm} \hat{X}_i = \{x_1,...,x_{n+1},C\}
$
\\ \vspace{4mm}
$
O[n] = |\hspace{1mm} s^+[X_i]|
$
\\ \vspace{2mm}
$
O[n+1] = \hat{O}[n] = |\hspace{1mm} s^+[\hat{X}_i]|
$
\end{center}
\hspace{4mm}
For general solutions $s^+$
\begin{center}
$
s^+[\hat{X}_i] \supseteq s^+[X_i] 
$
\\ \vspace{2mm}
$
\Rightarrow |s^+[\hat{X}_i]| \geq |s^+[X_i]|
$
\\ \vspace{2mm}
$
\therefore \hat{O}[n] = O[n+1] \geq O[n]
$
\end{center}










%\newpage
%\section{Convergent Complexity}
%\subsection{Definition}
%Define Convergent Complexity; the set of solutions with complexity satisfying
%\begin{center}
%$
%limit_{n \rightarrow \infty} \frac{O[n+1]}{O[n]} = c
%$
%\end{center}
%where c is a constant


%\subsection{Derivative Property of Convergent Solutions}
%There exists an nth derivative equal to zero
%\begin{center}
%$
%limit_{n \rightarrow \infty} \frac{O[n+1]}{O[n]} = c
%$
%\end{center}









%\subsection{Show c=1 for all convergent solutions}
%O[n] is an integer valued, non-decreasing function
%\begin{center}
%$
%O[n] \in \mathbb{Z}
%$
%\\ \vspace{2mm}
%$
%O[n+1] \geq O[n]; \hspace{2mm} f_{n+1}[n] = O[n+1] - O[n] \geq 0
%$
%\\ \vspace{6mm}
%$
%limit_{n \rightarrow \infty} \frac{O[n+1]}{O[n]} = limit_{n \rightarrow \infty} \frac{O[n] + f_{n+1}[n]}{O[n]}
%$
%\\ \vspace{2mm}
%$
%=  limit_{n \rightarrow \infty}(\frac{O[n]}{O[n]} +   \frac{f_{n+1}[n]}{O[n]})
%$
%\\ \vspace{2mm}
%$
%= 1 +   limit_{n \rightarrow \infty} \frac{f_{n+1}[n]}{O[n]} = c
%$
%\end{center}









%\subsection{Convergent Inductive Function}
%\begin{center}
%$
%limit_{n \rightarrow \infty} \frac{O[n+1]}{O[n]} = limit_{n \rightarrow \infty} \frac{O[n] + f_{n+1}[n]}{O[n]}
%$
%\\ \vspace{2mm}
%$
%=  limit_{n \rightarrow \infty}(\frac{O[n]}{O[n]} +   \frac{f_{n+1}[n]}{O[n]})
%$
%\\ \vspace{2mm}
%$
%= 1 +   limit_{n \rightarrow \infty} \frac{f_{n+1}[n]}{O[n]} = c
%$
%\\ \vspace{2mm}
%$
%= limit_{n \rightarrow \infty} \frac{f_{n+1}[n]}{O[n]} = c - 1
%$
%\end{center}


\section{Complexity of Canonical Instructions}
\begin{center}
$
c := a \leftarrow l[X_n]
$
\end{center}

\section{Complexity of Computational Operations}
\subsection{+}
Express the bounds of complexity for Computational Operation +



% Definition of Divergent Complexity
\section{Divergent Complexity}

\subsection{Definition of Converges to}
\begin{center}
\vspace{1mm}
$f[n]$ $converges$ $to$ $C$ = $convergent [f [n],C] = a_o; \hspace{1mm} a_o \in \{\mathbb{T},\mathbb{F}\} =$
\\ \vspace{6mm}
$
|C - f[n+1]| < |C - f[n]| \hspace{2mm} \forall n
$
\\ \vspace{2mm}
$
\land
$
\\ \vspace{2mm}
$
\not \exists K : | C - f[\hat{n}] | > K \hspace{4mm} \forall n; K > 0
$
\end{center}

\subsubsection{Notation}
C is commonly denoted by a limit
\begin{center}
$
C = lim_{n \rightarrow \infty} f[n]
$
\end{center}

\subsection{Definition of General Convergence}
\begin{center}
\vspace{1mm}
$f[n]$ is $convergent$ = $convergent \lbrack f \lbrack n \rbrack \rbrack = a_o; \hspace{1mm} a_o \in \{\mathbb{T},\mathbb{F}\} =$
\\ \vspace{6mm}
$
\exists C :
$
\\ \vspace{2mm}
$convergent[f[n],C] == \mathbb{T}$
\end{center}
Alternatively
\begin{center}
\vspace{1mm}
$f[n]$ is $convergent$ = $convergent \lbrack f \lbrack n \rbrack \rbrack = a_o \in \{\mathbb{T},\mathbb{F}\} =$
\\ \vspace{2mm}
$
\exists C :
$
\\ \vspace{2mm}
$f[n]$ $converges$ $to$ $C$
\end{center}

\subsection{Definition of Divergence}
\begin{center}
\vspace{1mm}
$
diverges \lbrack f \lbrack n \rbrack \rbrack = \lnot converges[f[n]] = d_o; \hspace{1mm} d_o \in \{\mathbb{T},\mathbb{F}\}
$
\\ \vspace{4mm}
$
:= \not \exists C : convergent[f[n],C] == \mathbb{T}
$
\end{center}





\subsection{Alternate Definition of Divergence}
\begin{center}
\vspace{1mm}
$diverges \lbrack f \lbrack n \rbrack \rbrack = \lnot converges[f[n]] = d_o; \hspace{1mm} d_o \in \{\mathbb{T},\mathbb{F}\}$
\\ \vspace{4mm}
$
= convergent[f[n],C] == \mathbb{F} \hspace{4mm} \forall C
$
\end{center}
\subsubsection{Proof of Equivalence; Alternate Definition of Divergence}





\subsection{?Necessary or Sufficient? Criteria 1 For Divergence}
? The derivative as a function of K ?
Function f[n] diverges if the $K^{th}$ derivative of f[n] is strictly increasing
\vspace{1mm}
\begin{center}
$
diverges \lbrack f \lbrack n \rbrack \rbrack := \not \exists C : convergent[f[n],C] == \mathbb{T}
$
\\ \vspace{4mm}
$
\Longleftrightarrow
$
\\ \vspace{4mm}
$
\Delta_n^{K+1} f[n] > \Delta_n^{K}f[n] \hspace{3mm} \forall K
$
\end{center}
Alternatively
\begin{center}
$
\Delta_n^{K+1} f[n] - \Delta_n^{K}f[n] > 0 \hspace{3mm} \forall K
$
\\ \vspace{4mm}
$
\Delta_{n}^{K+2} > 0 \hspace{2mm} \forall K
$
\end{center}






\subsubsection{Criteria 1; Proof of Necessity and Sufficiency}
\begin{center}
$
diverges \lbrack f \lbrack n \rbrack \rbrack = d_o; \hspace{1mm} d_o \in \{\mathbb{T},\mathbb{F}\}
$
\\ \vspace{4mm}
$
= \not \exists C : convergent[f[n],C] == \mathbb{T}
$
\end{center}
Let
\begin{center}
$
f[n] :
$
\\ \vspace{2mm}
$
\Delta_{n}^{K+2} > 0 \hspace{3mm} \forall K
$
\end{center}




\subsection{?Necessary or Sufficient? Criteria 2 For Divergence}
Function f[n] diverges if the Derivative as a function of K does not Converge
\begin{center}
\vspace{2mm}
$
f[n]\hspace{1mm} is \hspace{1mm} Divergent = Divergent[f[n]] = a_o \in \{\mathbb{T},\mathbb{F}\} :=
$
\\ \vspace{4mm}
$
\not \exists c: lim_{n\rightarrow \infty} \Delta_n^K f[n] = c
$
\end{center}








\subsubsection{Criteria 2; Proof of Necessity and Sufficiency}
\begin{center}
$
diverges \lbrack f \lbrack n \rbrack \rbrack = d_o; \hspace{1mm} d_o \in \{\mathbb{T},\mathbb{F}\}
$
\\ \vspace{4mm}
$
= \not \exists C : convergent[f[n],C] == \mathbb{T}
$
\end{center}






\subsection{Verbal Expressions}
\begin{center}
\vspace{1mm}
$
f[n]$ $diverges$ = $f[n]$ $is$ $divergent =
$
\\ \vspace{2mm}
$
f[n]$ $is$ $not$ $convergent= f[n]$ $does$ $not$ $converge
$
\end{center}

% Definition of Divergent Function
\subsection{Definition of Divergent Function}



\subsubsection{Definition 1}
Define Divergent Function f[n] having strictly increasing $K^{th}$ derivative
\vspace{1mm}
\begin{center}
$
f[n]\hspace{1mm} is \hspace{1mm} Divergent = Divergent[f[n]] = a_o \in \{\mathbb{T},\mathbb{F}\} :=
$
\\ \vspace{4mm}
$
\Delta_n^{K+1} f[n] > \Delta_n^{K}f[n] \hspace{3mm} \forall K
$
\end{center}
Alternatively
\begin{center}
$
\Delta_n^{K+1} f[n] - \Delta_n^{K}f[n] > 0 \hspace{3mm} \forall K
$
\\ \vspace{4mm}
$
\Delta_{n}^{K+2} > 0 \hspace{2mm} \forall K
$
\end{center}

\subsubsection{Definition 2}
\begin{center}
\vspace{2mm}
$
f[n]\hspace{1mm} is \hspace{1mm} Divergent = Divergent[f[n]] = a_o \in \{\mathbb{T},\mathbb{F}\} :=
$
\\ \vspace{4mm}
$
\not \exists c: lim_{n\rightarrow \infty} \Delta_n^K f[n] = c
$
\end{center}

\subsubsection{Proof of Equivalence Definition 1 $\Leftrightarrow$ Definition 2}

\subsubsection{Sufficient Proof}
\begin{center}
\vspace{2mm}
$
f[n]\hspace{1mm} is \hspace{1mm} Divergent \Longleftrightarrow
$
\\ \vspace{4mm}
$
\not \exists c: lim_{n\rightarrow \infty} \Delta_n^K f[n] = c
$
\end{center}



\subsubsection{Necessary Proof}
\begin{center}
\vspace{2mm}
$
f[n]\hspace{1mm} is \hspace{1mm} Divergent \Longleftrightarrow
$
\\ \vspace{4mm}
$
\not \exists c: lim_{n\rightarrow \infty} \Delta_n^K f[n] = c
$
\end{center}








Proof by contradiction of definition of limit
Using the definition of increasing convergence for a discrete function$^*$




\subsection{Defintion}
Decision problem $\hat{D}$ with solution $s^+$ has divergent total complexity $O[n]$ if
\begin{center}
$
lim_{n \rightarrow \infty} \frac{O[n+1]}{O[n]} \hspace{2mm} diverges
$
\end{center}







\subsection{Divergent Problems}
\begin{center}
$
\mathcal{\hat{D}} := \{ \hat{D}_1,\hat{D}_2,...\} :
$
\\ \vspace{2mm}
$
lim_{n \rightarrow \infty} \frac{O[n+1]}{O[n]} \hspace{2mm} diverges \hspace{3mm} \forall s^+ \in S^+_i, \hspace{1mm} \hat{D}_i \in \hat{\mathcal{D}}
$
\end{center}





%\subsection{Derivative Property of Divergent Solutions}
%\begin{center}
%$
%lim_{n \rightarrow \infty} \hspace{2mm} O[n+1] - O[n]$ diverges
%\end{center}

%\subsubsection{Proof}




%\subsection{Theorem of Divergent Subfunctions}
%If an (inductive) subfunction of $s^+$ diverges, the solution is divergent
%\begin{center}
%$
%f_{n+1} = \sum g_{n+1}
%$
%\\ \vspace{2mm}
%$
%limit \frac{g[n+1]}{g[n]} diverges
%$
%\\ \vspace{3mm}
%$
%\exists g_{n+1} : limit \frac{g_{n+1}}{O[n]} diverges \Longrightarrow limit \frac{O[n+1]}{O[n]} diverges
%$
%\end{center}

%\subsection{Proof}




\subsection{The Set of Polynomial Solutions and the Set of Divergent Solutions are disjoint}
\begin{center}
\vspace{2mm}
$
\mathbb{P} \cap \hat{D} = \emptyset
$
\end{center}

\subsection{Proof}
Let D $\in \hat{\mathcal{D}}$
\begin{center}
$
lim_{n \rightarrow \infty} \frac{O[n+1]}{O[n]} \hspace{2mm} diverges$ by definition
\end{center}
\vspace{4mm}
Assume D $\in \mathbb{P}$
\begin{center}
$
lim_{n \rightarrow \infty} \frac{O[n+1]}{O[n]} = 1
$
\\ \vspace{2mm}
$
 lim_{n \rightarrow \infty} \frac{O[n+1]}{O[n]}$ = 1 contradicts the definition of Divergent Problems
\\ \vspace{2mm}
$
\therefore D \in \hat{\mathcal{D}} \Rightarrow D \not \in \mathbb{P}
$
\end{center}
\vspace{12mm}
Let D $\in \mathbb{P}$
\begin{center}
$
lim_{n \rightarrow \infty} \frac{O[n+1]}{O[n]} = 1$ by property of Polynomial complexity
\end{center}
\vspace{4mm}
Assume D $\in \hat{D}$
\begin{center}
$
lim_{n \rightarrow \infty} \frac{O[n+1]}{O[n]}$ diverges
\\ \vspace{2mm}
$
 lim_{n \rightarrow \infty} \frac{O[n+1]}{O[n]}$ diverges contradicts a property of Polynomial complexity
\\ \vspace{2mm}
$
\therefore D \in \mathbb{P} \Rightarrow D \not \in \hat{\mathcal{D}}
$
\\ \vspace{2mm}
$
\therefore \mathbb{P} \cap \hat{\mathcal{D}} = \emptyset
$
\end{center}






% Definition of Workers
\newpage
\section{Workers}

\section{Definition of Parallel Programs}

\section{Can Parallel Workers transform a Divergent task into a Polynomial Task}


\section{Conservation of Computation}









% N Sum Problem
\newpage
\section{Sum to N Problem with 2 integers}
% Zero Order Space Complexity Solution
\subsection{State formal definition of Sum to N : $x_i + x_j == N$}
\begin{center}
\vspace{1.5mm}
$
X_n= \{x_1,...,x_n\}
$
\\ \vspace{6mm}
$
D := f \lbrack X_i,N \rbrack = a_{o} \in \{\mathbb{T}, \mathbb{F}\} \hspace{3mm} \forall X_i
$
\\ \vspace{6mm}
$
s^+[X_n]= P[X_n] :
$
\\ \vspace{2mm}
$
(P \lbrack X_i \rbrack = y_{o} == a_{o} \hspace{3mm} \forall X_i) \hspace{2mm} \cap \hspace{2mm} (P[{X}_{n+1}] \supseteq P[X_n] \hspace{3mm} \forall X_{n+1})
$
\\ \vspace{6mm}
$
s^+ = \{ s_1,s_2,...,s_{O_T \lbrack n \rbrack }, b_1, b_2,...,b_{O_S \lbrack n \rbrack},y_o \} = \{ \mathcal{L},\mathcal{M},y_o\}
$
\\ \vspace{6mm}
$
D = f[X_i] = \hspace{2mm} \exists x_j,x_k \in X_n \hspace{4mm} j\neq k: 
$
\\ \vspace{2mm}
$
x_j + x_k == N
$
\end{center}





\subsection{Express a formal solution : $O_S[n] \sim n^0$}
\begin{center}
\vspace{1.5mm}
$
s^+ = \{ s_1,s_2,...,s_{O_T \lbrack n \rbrack }, b_1, b_2,...,b_{O_S \lbrack n \rbrack},y_o \} = \{ \mathcal{L},\mathcal{M},y_o\}
$
\\ \vspace{2mm}
$
s_1 = y_o \leftarrow \mathbb{F};
$
\\ \vspace{2mm}
$
\forall i < n \hspace{2mm}, \hspace{2mm} n \geq j > i 
$
\\ \vspace{2mm}
$
s_2,s_3,s_8,s_9,...,s_{3ij-4},s_{3ij-3}...,s_{3n(n-1)-4},s_{3n(n-1)-3} =  b_1 \leftarrow x_i + x_j 
$
\\ \vspace{2mm}
$
s_4,s_5,s_{10},s_{11},...,s_{3ij-2},s_{3ij-1}...,s_{3n(n-1)-2},s_{3n(n-1)-1} = b_1 \leftarrow b_1 == N
$
\\ \vspace{2mm}
$
s_6,s_7,s_{12},s_{13}...,s_{3ij},s_{3ij+1}...,s_{3n(n-1)},s_{3n(n-1)+1} = y_o \leftarrow y_o \lor b_1
$
\\ \vspace{2mm}
$
s^+ = \{y_o \leftarrow \mathbb{F},y_o \leftarrow y_o \lor (x_i + x_j == N) \hspace{3mm} \forall i,j > i \hspace{1mm}| \hspace{1mm} b_1,y_o\}
$
\end{center}





\subsection{Prove $s^+$ satisfies the subfunction condition of solutions: $P[X_{n+1}] \supseteq P[X_n] \hspace{3mm} \forall X_{n+1}$}
\begin{center}
\vspace{2mm}
$
X_n= \{x_1,x_2,...,x_n\}; \hspace{2mm} X_{n+1} = \{x_1,x_2,...,x_n,x_{n+1}\}
$
\\ \vspace{2mm}
$
s^+ = \{ s_1,s_2,...,s_{O_T \lbrack n \rbrack }, b_1, b_2,...,b_{O_S \lbrack n \rbrack},y_o \} = \{ \mathcal{L},\mathcal{M},y_o\}
$
\\ \vspace{2mm}
$
s_{n+1}^+ = s^+ \cup \hat{s}^+ 
$
\\ \vspace{2mm}
$
s_1 = y_o \leftarrow \mathbb{F};
$
\\ \vspace{2mm}
$
\forall i < n \hspace{2mm}, \hspace{2mm} n \geq j > i 
$
\\ \vspace{2mm}
$
s_2,s_3,s_8,s_9,...,s_{3ij-4},s_{3ij-3}...,s_{3n(n-1)-4},s_{3n(n-1)-3} =  b_1 \leftarrow x_i + x_j 
$
\\ \vspace{2mm}
$
s_4,s_5,s_{10},s_{11},...,s_{3ij-2},s_{3ij-1}...,s_{3n(n-1)-2},s_{3n(n-1)-1} = b_1 \leftarrow b_1 == N
$
\\ \vspace{2mm}
$
s_6,s_7,s_{12},s_{13}...,s_{3ij},s_{3ij+1}...,s_{3n(n-1)},s_{3n(n-1)+1} = y_o \leftarrow y_o \lor b_1
$
\\ \vspace{6mm}
$
\forall k < n+1
$
\\ \vspace{2mm}
$
s... = b_1 \leftarrow x_k + x_{n+1}
$
\\ \vspace{2mm}
$
s... = b_1 \leftarrow b_1 == N
$
\\ \vspace{2mm}
$
s... = y_o \leftarrow y_o \lor b_1
$
\\ \vspace{8mm}
$
s^+ = \{y_o \leftarrow \mathbb{F},y_o \leftarrow y_o \lor (x_i + x_j == N) \hspace{3mm} \forall i,j > i \hspace{1mm}| \hspace{1mm} b_1,y_o\}
$
\\ \vspace{4mm}
$
\hat{s}^+ = \{y_o \leftarrow y_o \lor (x_k + x_{n+1} == N) \hspace{2mm} \forall k < n+1 | \hspace{1mm} b_1,y_o \}
$
\\ \vspace{4mm}
$
s^+_{n+1} = \{y_o \leftarrow \mathbb{F},y_o \leftarrow y_o \lor (x_i + x_j == N) \hspace{3mm} \forall i,j > i \hspace{1mm}| \hspace{1mm} b_1,y_o\} \hspace{2mm}\cup 
$
\\ \vspace{3mm}
$
\{y_o \leftarrow y_o \lor (x_k + x_{n+1} == N) \hspace{2mm} \forall k < n+1 | \hspace{1mm} b_1,y_o \}
$
\\ \vspace{6mm}
$
s_{n+1}^+ = s^+ \cup \hat{s}^+ = P[X_{n+1}] \supseteq P[X_n] = s^+
$
\end{center}






\subsection{Determine $O[n], O_S[n], O_T[n],f_{n+1}[n],f^T_{n+1}[n],f^S_{n+1}[n]$ for the above solution}
\begin{center}
$
O_S[n] = |y_o| + |b_1| = 2
$
\\ \vspace{2mm}
$
O_T[n] = 3n(n-1) + 1 = 3n(n-1) -1 + O_S[n]
$
\\ \vspace{2mm}
$
O[n] = 3n(n-1) + 3 = 3n^2 -3n + 3
$
\\ \vspace{4mm}
$
f^S_{n+1}[n] = 0
$
\\ \vspace{2mm}
$
f^T_{n+1}[n] = 6n
$
\\ \vspace{2mm}
$
f_{n+1}[n] = f^S_{n+1}[n] + f^T_{n+1}[n]
$
\end{center}


\subsection{Verify $O[n+1]$ = $O[n]$ + $f_{n+1}[n]$}
\begin{center}
$
O[n+1] = O[n] + \hat{O}[n]
$
\\ \vspace{2mm}
$
3(n+1)^2 -3(n+1) + 3 =  3n^2 -3n + 3 + 6n
$
\\ \vspace{2mm}
$
3n^2 + 6n + 3 - 3n -3 + 3 = 3n^2 + 3n + 3
$
\\ \vspace{2mm}
$
3n^2 + 3n+ 3 = 3n^2 + 3n + 3
$
\end{center}







\subsection{Show $s^+$ has Polynomial Complexity by the definition of Total Polynomial Complexity}
\begin{center}
\vspace{2mm}
$
O[n] = 3n^2 -3n + 3
$
\end{center}

\subsection{Show the limit $_{n \rightarrow \infty}\frac{O[n+1]}{O[n]}$ does not Diverge}
\begin{center}
$
limit_{n \rightarrow \infty}\frac{O[n+1]}{O[n]} =
$
\\ \vspace{2mm}
$
limit_{n \rightarrow \infty}\frac{3n^2 + 3n + 3}{3n^2 - 3n + 3} =
$
\\ \vspace{2mm}
$
limit_{n \rightarrow \infty}(\frac{3n^2 - 3n + 3}{3n^2 - 3n + 3} + \frac{6n}{3n^2 - 3n + 3}) =
$
\\ \vspace{2mm}
$
limit_{n \rightarrow \infty}(1 + \frac{6n }{3n^2 - 3n + 3}) = 1
$
\end{center}





% Knapsack Problem
\newpage
\section{The Knapsack Problem}


% Knapsack Problem Verbal Definition
\subsection{The Knapsack Problem}
The Knapsack Problem is a famous problem in computer science which asks if objects can be stored in a knapsack. Typically the problem is designed
with two constraints, weight and value. Given objects $x_i$, each with a respective weight $w_i$ and value $v_i$, does there exist a combination of
objects lighter than input weight $W$ and more valuable than input value $V$?


% Knapsack Problem Formal Defintion
\subsection{Formal Definition}
\begin{center}
\vspace{8mm}
$
X_n = \{x_1,x_2,...,x_n\} = \{ \{w_1,v_1\},\{w_2,v_2\},...,\{w_n,v_n\} \}
$
\\ \vspace{8mm}
$
I = \{i_1,i_2,...,i_n\}: i_l \in \{0,1\} \hspace{2mm} \forall i_l \in I
$
\\ \vspace{8mm}
$
D := f[X_n,W,V] = a_o \in \{\mathbb{T},\mathbb{F}\} = \exists I:
$
\\ \vspace{2mm}
$
(\sum_{j=1}^n i_j w_j < W) \land (\sum_{j=1}^n i_j v_j \geq V)
$
\end{center}
\vspace{3mm}








\subsection{Solve for C[n]}

\subsubsection{Expressing $I$ as a binary number}
\vspace{4mm}
\begin{center}
$
I = \{i_1,i_2,...,i_n\}: i_l \in \{0,1\} \hspace{2mm} \forall i_l \in I
$
\end{center}
Valid combinations of I
\begin{center}
$
I_{valid} =
$
\\ \vspace{2mm}
$
\{ \{0,0,0,...,0,0,1\}, \{0,0,0,...,0,1,0\}, \{0,0,0,...,0,1,1\}, ... ,\{1,1,1,...,1,1,1\}\}
$
\\ \vspace{4mm}
$
C[n] = |I_{valid}[n]| = 2^n - 1
$
\end{center}



\subsubsection{Using a sum of combinations of inputs $x_i$}
\vspace{4mm}
\begin{center}
$
X_n = \{x_1,x_2,...,x_n\} = \{ \{w_1,v_1\},\{w_2,v_2\},...,\{w_n,v_n\} \}
$
\end{center}
Valid combinations of $x_i$
\begin{center}
$
X_{valid}[n] =
$
\\ \vspace{2mm}
$
\{x_1\} \cup \{x_2\} \cup ... \cup \{x_n\} \cup \{x_1,x_2\}  \cup \{x_1,x_3\} \cup ... \cup \{x_{n-1},x_{n}\} \cup ... \cup \{x_1,x_2,...,x_n\}
$
\\ \vspace{2mm}
$
= \ _{X_n}C_1 \cup \hspace{1mm} _{X_n}C_2 \cup ... \cup \hspace{1mm}_{X_n}C_n
$
\\ \vspace{6mm}
$
C[n] = |X_{valid}[n]| = {\sum _{j=1}^{n}}\hspace{2mm}{_{n}C{_j}}
$
\end{center}




\subsubsection{Verify consistency}
\begin{center}
\vspace{4mm}
$
C[n] = |X_{valid}[n]| = |I_{valid}[n]|
$
\\ \vspace{4mm}
$
= 2^n - 1 =  {\sum _{j=1}^{n}}\hspace{2mm}{_{n}C{_j}} = {_{n}C{_1}} + {_{n}C{_2}} + ... + {_{n}C{_n}}
$
\\ \vspace{4mm}
$
= 2^n -1 = 2^n -1
$
\end{center}





\subsection{Express a solution $s^+$ of the Knapsack Problem}
\subsection{Prove $s^+$ satisfies the subfunction condition of solutions}
\subsection{Determine $O[n],O_T[n],O_S[n],f_{n+1}[n]$}
\subsection{Show $s^+$ $\notin$ $\mathbb{P}$}
\subsection{Express the Solution Space $\mathbb{S}$ for The Knapsack Problem}
\subsection{Prove a lower bound for all solutions $s^+ \in S^+ := O_{lower}[n]$}
\subsection{Prove $D$ has Divergent Complexity}




\newpage
\section*{Appendix}



\section{Criticism of Overloaded Equivalence}
In computer science, it is convention to overload equivalence $=$
\begin{center}
$
a_i = a_i \hspace{3mm} \forall a_i \in \Omega
$
\end{center}
\vspace{3mm}
Consider standard C++ syntax\\
int x = 3;\\
int y = 4;\\
int z = x + y;\\ \\
Int x is not inherently equal to 3. Rather, we are creating an open space "x" for a value and setting the value to 3. Similarly, z is not inherently equal
to the value of x + y. Rather, we are creating an open space "z" for a value and setting the value to the sum of x and y which have already been set.
\begin{center}
$
x \leftarrow 3
$
\\ \vspace{2mm}
$
y \leftarrow 4
$
\\ \vspace{2mm}
$
z \leftarrow x + y
$
\end{center}

\section{Existence of $O_{max}[n]$}
\subsection{Proof}


\subsubsection{Alternate Definition; Left Hand Derivative}
Some sources define
\begin{center}
$
\Delta_n^1 f[n] = f[n] - f[n-1]
$
\end{center}


\newpage
\section*{Citations}
\lbrack1\rbrack \hspace{1mm} $https://kapeli.com/cheat\_sheets/LaTeX\_Math\_Symbols.docset/Contents/Resources/Documents/index$\\
\lbrack2\rbrack \hspace{1mm} $chat.openai.com$\\
\lbrack3\rbrack \hspace{1mm} $google.com$\\
\lbrack4\rbrack \hspace{1mm} $wikipedia.org$\\
\lbrack3\rbrack \hspace{1mm} $wolframalpha.com$\\
\lbrack4\rbrack \hspace{1mm} $youtube.com$\\



\end{document}